\section{The protocol}
\label{sec:protocol}

A certificate is only as meaningful as the object it is applied to. This
section describes the three decisions that turn the machinery of
Section~\ref{sec:cert} into a procedure that produces an honest number for a
trained network. Figure~\ref{fig:pipeline} shows the pipeline; each subsection
below justifies one arrow.

\begin{figure}[tb]
\centering
\begin{tikzpicture}[font=\small,
  box/.style={draw, rounded corners=2pt, align=center, inner sep=3pt,
              minimum height=14mm, text width=1.5cm, fill=white},
  arr/.style={-{Stealth[length=1.8mm]}, semithick},
  lab/.style={midway, above=0.9cm, font=\scriptsize}]
\node[box] (net)  {network output\\ $\psi_\theta$};
\node[box, right=1.05cm of net] (fil)
      {guided DST filter\\ $(K_x,K_y)$};
\node[box, right=1.05cm of fil] (cnt)
      {verified counts\\ $N(a)$, $N(b)$};
\node[box, right=1.05cm of cnt] (enc)
      {Temple\\ bound $I$};
\node[box, right=1.05cm of enc] (trn)
      {certified angle\\ for $\psi_\theta$};
\draw[arr] (net) -- (fil) node[lab] {$\psi_K$};
\draw[arr] (fil) -- (cnt) node[lab] {$\eta,\rho$};
\draw[arr] (cnt) -- (enc) node[lab] {count $1$};
\draw[arr] (enc) -- (trn) node[lab] {Lemma~\ref{lem:transport}};
\draw[arr, dashed] (enc.south) -- ++(0,-1.0) -| (fil.south)
      node[pos=0.25, below, font=\scriptsize] {re-filter if worse};
\end{tikzpicture}
\caption{The certification pipeline. The counters are called for each
candidate truncation pair; the filter keeps the narrowest enclosure among
those whose verified state index agrees with that of the unfiltered output
(Section~\ref{sec:filter}). The parenthesised arrows are the feedback that
makes the filter guided rather than fixed.}
\label{fig:pipeline}
\end{figure}

\subsection{A residual whose energy cannot leak}
\label{sec:leakage}

The training objective is the squared residual of \eqref{eq:residual}, a term
that fixes the scale, a boundary term and an energy term,
\begin{equation}
  \begin{split}
  \mathcal L(\theta) &=
  \operatorname*{mean}_{i}\big(\Hh\psi_\theta-\eta(\psi_\theta)\psi_\theta\big)_i^2
  + w_n\big(\langle\psi_\theta,\psi_\theta\rangle_h-1\big)^2 \\
  &\quad + w_b\big(\psi_\theta(a)^2+\psi_\theta(b)^2\big)
  + w_e\,\eta(\psi_\theta),
  \end{split}
  \label{eq:loss}
\end{equation}
with $(w_n,w_b,w_e)=(100,100,10)$ in one dimension. In two dimensions the
residual and the energy are evaluated on the unit iterate
$u=\psi_\theta/\norm{\psi_\theta}_h$, the Dirichlet condition is imposed
exactly by replacing the network output with
$\psi_\theta\leftarrow M\odot\psi_\theta$, where
$M(x,y)=(x-a)(b-x)(y-a)(b-y)$, so the boundary term vanishes identically and
$w_e=10$ is kept. In both dimensions $\eta(\psi_\theta)$ is the Rayleigh
quotient of the iterate, treated as a constant in the backward pass (it is the
minimizer over $E$, so detaching it leaves the gradient correct).

Four properties of \eqref{eq:loss} matter for certification. First, because
the residual uses the same $\eta$ as the certificate, the training objective
and the certified quantity coincide; a network trained on a different energy
(a learnable scalar, or a reference value) is optimizing a different residual,
and the certificate then measures that other objective. Second, the scale must
be fixed by an explicit term: without one $\psi\equiv0$ is an exact stationary
point of the residual, the loss is quadratic in the amplitude, and the amplitude is
left wherever the initialization put it---in our unnormalized two-dimensional
run it stayed as small as $\norm{\psi}^2=6\times10^{-5}$ while the shape had
already converged. One dimension fixes the scale with the penalty
$(\norm{\psi}_h^2-1)^2$; two dimensions divide by the norm inside the loss
instead. Third, the boundary condition: in one dimension the soft penalty of
weight $100$ leaves endpoint values of order $10^{-4}$ and contributes a
negligible part of the residual, whereas in two dimensions the same penalty is
diluted by the $O(1/N)$ fraction of boundary nodes and leaves the residual at
$O(1)$ (\ref{app:impl} records both failures), which is why the mask is used
there. Fourth, the energy term $w_e\eta$ is variational: it pulls the iterate
toward the lowest Rayleigh quotient and selects the ground state within the
near-null space of the residual; because $\eta$ is computed from the iterate,
it carries no external value.

\subsection{Certificate-guided spectral filtering}
\label{sec:filter}

A trained network contains high-frequency components that contribute almost
nothing to the Rayleigh quotient---which is stationary to first order---but
dominate the residual, which is not. Filtering them out turns an
uncertifiable output into a certifiable one. Under Dirichlet conditions the
natural basis is the tensor product of sine functions,
\begin{equation}
  \psi(x_i,y_j)\approx\sum_{p=1}^{K_x}\sum_{q=1}^{K_y} c_{pq}
  \sin\frac{p\pi x_i}{L}\sin\frac{q\pi y_j}{L},
  \label{eq:dst}
\end{equation}
with coefficients obtained by a two-dimensional type-I discrete sine
transform (DST-I) \citep{boyd2001}, which is separable, orthogonal up to a
diagonal scaling, and self-inverse.

The truncation order is chosen without any reference solution. For each
candidate pair $(K_x,K_y)$ on a prescribed schedule we filter, certify, and
record $(\text{width}, \text{state index})$. The selected pair is the one with
the smallest certified width \emph{among those whose certified state index
equals that of the unfiltered input}. The state-consistency constraint is not
optional: without it the filter happily projects an excited state onto the
ground-state sine, producing a much narrower certificate for the wrong
function. On the runs below the filter changes $\psi$ by a relative $L^2$
amount between $4\times10^{-6}$ and $2\times10^{-3}$, and reduces the residual
norm by a factor between $1.2$ and $7\times10^{9}$---below threefold in half
of the eighteen runs of the main sweep.

\subsection{Transporting the bound to the raw output}
\label{sec:transport}

The certificate of Section~\ref{sec:filter} refers to $\psi_K$, not to the
network output $\psi$. The gap can be bounded by an elementary inequality, so
that the network output itself receives a certified angle.

\begin{lemma}[Transport]\label{lem:transport}
Let $u=\psi/\norm{\psi}_h$, $u_K=\psi_K/\norm{\psi_K}_h$, and let $\varphi$ be
a unit eigenfunction. Then
\begin{equation}
  \sin\angle(u,\varphi)\ \le\ 2\norm{u-u_K}_h+\sin\angle(u_K,\varphi).
  \label{eq:transport}
\end{equation}
\end{lemma}
\begin{proof}
Since $\sin\angle(u,\varphi)=\norm{u-(u\cdot\varphi)\varphi}_h$ and
$\abs{(u_K-u)\cdot\varphi}\le\norm{u-u_K}_h$,
\begin{align*}
\sin\angle(u,\varphi)
&\le \norm{u-u_K}_h+\norm{u_K-(u_K\!\cdot\!\varphi)\varphi}_h
   +\abs{(u_K-u)\cdot\varphi}\\
&\le 2\norm{u-u_K}_h+\sin\angle(u_K,\varphi). \qedhere
\end{align*}
\end{proof}

The first term is computed from the two saved functions; the second is
Theorem~\ref{thm:dk} applied to $\psi_K$. The implementation uses the
conservative factor $2$ of \eqref{eq:transport}; a factor $1$ also follows from
the Lipschitz property of $\sin$ on $[0,\pi/2]$. The inequality is what allows
the tables to report a bound for the network output even when the unfiltered
output is not certifiable, and it answers the question "what exactly is being
certified?" with a number rather than a caveat.

\subsection{What a row of the results table contains}
\label{sec:row}

For each run the protocol reports: the certified state index and whether it
matches the target; the energy $\eta$ and the distance to an independently
computed $\lambda_h$; the residual before and after filtering; the certified
width before and after filtering; the certified angle bound, the transported
bound for the raw output, and the true discrete angle (computed only for
validation, never used by the protocol); the relative $L^2$ change caused by
the filter; and whether the enclosure contains the reference. Reporting the
raw and filtered columns side by side is essential: the first is what the
network guarantees by itself, the second is what can be certified in its
neighbourhood. Reporting only the second would overstate the network;
reporting only the first would understate what the protocol can do.

The one case in which this row does not exist is the degenerate level: there
is no single state to describe, and the certificate of
Section~\ref{sec:cert} refuses. Its replacement---a record for the invariant
subspace, with the multiplicity measured by counting---is the subject of
Section~\ref{sec:subspace}, and it is the object that the experiments of
Section~\ref{sec:deg-results} report.
